\documentclass[11pt]{article}
\usepackage{mathblatt}
% Präambel des Setzers (werkzeuge/setzer-praeambel.tex, wortgleich); dazu \gteil (Muster T2U)
\makeatletter
\setlength{\mbpfbe}{0mm}% keine Punkte (pruefung.md § 5: nur beim Prüfstein)
\newlength{\gEin}\setlength{\gEin}{0pt}
\let\gSchrift\relax
\renewcommand{\pfkreuzzeile}[1]{\par\noindent\hangindent3.2em\hangafter1\hspace*{1.6em}$\square$\ #1\par}
\newcommand{\gkreuz}[1]{$\square$\ #1\hspace{2em}}
\newcommand{\gkreuzzeile}[2]{\par\noindent\hangindent2.2em\hangafter1\mbox{#1\,$\square$}\ #2\par}% Buchstabe, Option: eine je Zeile (neben Skizze)
% links, was man ansieht, rechts, was man tut; Spaltengrenze überall an derselben Stelle (gemeinsam.md „Satz“)
% v0.17: Fußeintrag „Nr: Wert“ – kurz in einer Box (nie getrennt), lang umbrechbar mit der Nummer am Anfang
% (vorher \mbox: lange Lösungen liefen über den Rand, AS40 138 pt); robust, weil \fusshilfe in die .aux schreibt
\NewDocumentCommand{\gfuss}{m m}{\sbox0{#1:~#2}\ifdim\wd0<0.45\textwidth\mbox{#1:~#2}\else\mbox{#1:}\nobreak\ #2\fi}
\newsavebox{\gbox}\newlength{\gLinks}
% v0.10 (Lehrer 13.10.): keine starre Halbierung – die linke Spalte ist so breit wie die Figur (+ 6 mm), die
% Antwort steht direkt daneben; ist die Figur breiter als 0,62 Zeile, steht die Antwort darunter
\newcommand{\gzwei}[2]{\par\smallskip\noindent\sbox{\gbox}{#1}%
  \setlength{\gLinks}{\dimexpr\wd\gbox+6mm\relax}%
  \ifdim\gLinks>0.62\linewidth
    \usebox{\gbox}\par\smallskip\noindent\begin{minipage}{\linewidth}#2\end{minipage}%
  \else
    \begin{minipage}[t]{\gLinks}\vspace{0pt}\usebox{\gbox}\end{minipage}%
    \begin{minipage}[t]{\dimexpr\linewidth-\gLinks\relax}\vspace{0pt}\raggedright #2\end{minipage}%
  \fi\par}
% 16.10. (Lehrer, PA10 S. 1): Text links, Koordinatensystem rechts daneben; ist das System breiter als 0,58 Zeile,
% steht es wie bisher unter dem Text
% Passt das System knapp nicht mehr auf die Seite (bis 25 % zu hoch), wird es auf den Rest verkleinert statt
% die Seite halb leer zu lassen
\newlength{\gRest}\newlength{\gHoch}
\newcommand{\gtextks}[2]{\sbox{\gbox}{#1}%
  \setlength{\gHoch}{\dimexpr\ht\gbox+\dp\gbox\relax}%
  \ifdim\pagegoal<\maxdimen\setlength{\gRest}{\dimexpr\pagegoal-\pagetotal-4\baselineskip\relax}%
    \ifdim\gHoch>\gRest\ifdim\gRest>0.75\gHoch
      \sbox{\gbox}{\resizebox*{!}{\gRest}{\usebox{\gbox}}}\fi\fi\fi
  \ifdim\wd\gbox>0.58\linewidth #2\par\smallskip\noindent\gpassend{\usebox{\gbox}}\par
  \else\begin{minipage}[t]{\dimexpr\linewidth-\wd\gbox-5mm\relax}\raggedright #2\end{minipage}\hfill
    \begin{minipage}[t]{\wd\gbox}\vspace{0pt}\usebox{\gbox}\end{minipage}\par\fi}
% v0.9 (12.10., pruefung.md § 5 „Marke mit B, Bedingungen rechts“): drittes Argument von pfaufg = Bedingungen
% der Prüfung (v0.10: \pfohnetr = TR durchgestrichen, „FOR“), klein grau im rechten Rand auf der Grundlinie der Nummer, Zeilen mit \\
\newcommand{\gbedrand}[2]{\ifblank{#2}{}{\rlap{\hspace*{\dimexpr#1+2mm\relax}%
  \parbox[t]{15mm}{\raggedright\scriptsize\leavevmode\color{mbgrau}\strut#2}}}}
% Teil eines Bündels (a, b …): Marke links im Rand, Bedingungen rechts im Rand, auf der Zeile des Teils
\newcommand{\gteilrand}[2]{\llap{\mbpfrandmarke{#1}\hspace*{\mbpfnr}}\gbedrand{\linewidth}{#2}}
% Teil eines Bündels: Buchstabe in eigener Spalte, Inhalt eingerückt (\gEin = Einzug, damit die Spaltengrenze von
% \gzwei an derselben Stelle bleibt)
\newenvironment{gbteil}[3]{\par\smallskip\noindent\gteilrand{#2}{#3}\makebox[\mbpfteil][l]{\textbf{#1}}%
  \begin{minipage}[t]{\dimexpr\linewidth-\mbpfteil\relax}\setlength{\gEin}{\mbpfteil}\raggedright\strut\ignorespaces}%
 {\end{minipage}\par\addvspace{6pt}}% v0.18 (PG80 Nr. 1): Abstand, damit die Linie von a) nicht an b) rutscht
% Grafik höchstens so breit wie die Spalte (Bündel nebeneinander, Koordinatensystem ohne Karo)
\newsavebox{\gpbox}
\newcommand{\gpassend}[1]{\sbox{\gpbox}{#1}\ifdim\wd\gpbox>\linewidth\typeout{GPASSEND: verkleinert \the\wd\gpbox\space auf \the\linewidth}\resizebox{\linewidth}{!}{\usebox{\gpbox}}\else\usebox{\gpbox}\fi}
% Skizzenfeld (v0.10, Lehrer 13.10.): Karo ab dem linken Rand über die Zeile, „Skizze“ klein grau oben links im Feld
% (wie \pfskizzenfeld der Vorlage); der rechte Rand bleibt den Bedingungen
\newcommand{\gskizze}[1]{\par\smallskip\noindent\begin{tikzpicture}
  \clip (0,0) rectangle ({\linewidth-1mm},{#1*10mm});
  \draw[black!16,line width=0.3pt,step=5mm] (0,0) grid ({\linewidth},{#1*10mm});
  \node[anchor=north west,inner sep=2pt,font=\scriptsize,text=mbgrau,fill=white] at (0,{#1*10mm}) {Skizze};
  \end{tikzpicture}\par}
\RenewDocumentEnvironment{pfaufg}{m m m}{%
  \begin{lrbox}{\mbaufgabebox}\begin{minipage}[t]{\linewidth}%
  \gSchrift\noindent\hspace*{\gEin}\mbpfrandmarke{#1}\gbedrand{\linewidth-\gEin}{#3}%
  \begin{minipage}[t]{\mbpfnr}\strut\textbf{#2}\end{minipage}%
  \begin{minipage}[t]{\dimexpr\linewidth-\gEin-\mbpfnr\relax}\raggedright\strut\ignorespaces}%
 {\end{minipage}%
  \end{minipage}\end{lrbox}%
  \setlength{\mbaufgabehoehe}{\dimexpr\ht\mbaufgabebox+\dp\mbaufgabebox\relax}%
  \ifdefempty{\gkopfnach}{\par\addvspace{5pt}\Needspace*{\mbaufgabehoehe}}{%
    % Stufenkopf, der mit dieser (ersten) Aufgabe zusammenbleibt: Platz für beide verlangen, dann den Kopf setzen
    % (12.10., Quadratische: ein Kopf stand allein unten, die Aufgabe mit Koordinatensystem auf der nächsten Seite)
    % v0.17 (AS40 „Terme deuten“ allein unten): hinter dem Kopf steht das Sprungziel (\special, ein Knoten, nach
    % dem die Seite brechen durfte) – darum \nobreak nach dem Kopf; Vorhalt 56 pt = Kopf mit Linie und Abständen
    \par\Needspace*{\dimexpr\mbaufgabehoehe+56pt\relax}\gkopfnach\par\nobreak\gdef\gkopfnach{}\addvspace{5pt}}%
  \noindent\usebox{\mbaufgabebox}\par\addvspace{5pt}}
% v0.18 (plan.md § 10, 15.10. l): Abschnittsbuchstabe A, B, C vor dem Namen im Kopf – als Vorsatz, damit der
% Name im TeX gleich bleibt (Bestellkarte, Sprungziele lesen ihn)
\def\gbstakt{}\newcommand{\gbstvor}[1]{\gdef\gbstakt{#1}}
\NewCommandCopy\gpfstufekopfalt\pfstufekopf
\RenewDocumentCommand{\pfstufekopf}{o m m}{\IfValueTF{#1}{\gpfstufekopfalt[#1]}{\gpfstufekopfalt}%
  {\ifx\gbstakt\empty\else\gbstakt\enspace\fi #2}{#3}\gdef\gbstakt{}}
% v0.18 (16.10. e): Begriff grau, ohne Nummer, ohne Rahmen über der ersten Aufgabe der Stufe
\newcommand{\gbegriff}[1]{\par\nobreak\noindent\hspace*{\mbpfnr}\parbox[t]{\dimexpr\linewidth-\mbpfnr\relax}{\raggedright\small\color{mbgrau}#1}\par\nobreak\addvspace{3pt}}
% v0.18 (16.10. g): Lösungsblatt – Kopf mit kräftigem Strich, Abschnitt klein grau als Zwischenzeile
\newcommand{\glkopf}[1]{\noindent{\large\bfseries Lösungen \textperiodcentered{} #1}\par\vspace{3pt}\hrule height 1.2pt\par\medskip}
\newcommand{\glzwischen}[1]{\par\addvspace{8pt}\Needspace*{4\baselineskip}\noindent{\footnotesize\color{mbgrau}#1}\par\nobreak\vspace{1pt}}
% v0.18 (15.10. o, w): Blattende – Kern ohne Rahmen, dann graue Zeile „Mehr: n Aufgaben“
\newcommand{\gkern}[1]{\par\addvspace{14pt}\Needspace*{6\baselineskip}\noindent{\bfseries Kern}\par\nobreak\smallskip
  {\small\raggedright #1\par}}
\newcommand{\gmehr}[1]{\par\addvspace{10pt}\noindent{\footnotesize\color{mbgrau}Mehr: #1 Aufgaben}\par}
% Stufenkopf vormerken: \gkopfvor{<Kopf>} – die nächste pfaufg setzt ihn, sobald beide zusammen Platz haben
\def\gkopfnach{}
\newcommand{\gkopfvor}[1]{\gdef\gkopfnach{#1}}
% Teilaufgabe mit Skizze links in gzwei: Buchstabe in eigener Spalte, darunter Text und Skizze
\newcommand{\gteil}[2]{\makebox[\mbpfteil][l]{\textbf{#1}}\parbox[t]{\dimexpr0.5\textwidth-\gEin-\mbpfnr-5mm-\mbpfteil\relax}{\raggedright #2}}
% Stufenkopf der Vorlage (mathblatt.sty 2026-10-12: Linie, Abstand, 13 pt, \pfjahre) ohne ihren Vorhalt von
% 14 Zeilen: der Setzer fordert vor dem Kopf 0,16 Seitenhöhe (Kopf + erste Aufgabe), mehr nicht
\setlength{\mbpfstufevorhalt}{0pt}
\makeatother
\newcommand{\karo}[1]{\par\smallskip\noindent\begin{tikzpicture}
  \clip (0,0) rectangle ({\linewidth-1mm},{#1*10mm});
  \draw[black!16,line width=0.3pt,step=5mm] (0,0) grid ({\linewidth},{#1*10mm});
  \end{tikzpicture}\par}
\newcommand{\kk}[1]{$\square$\,#1\hspace{0.9em}}
\newcommand{\linie}[1]{\,{\color{black!45}\rule[-1pt]{#1}{0.4pt}}\,}
\newcommand{\pkt}{\textperiodcentered{}}
% v0.14 (14.10.): Stufenkopf mit fünf Punkten (Option = Jahre für die Bestellkarte) oder „zuletzt 19“
\NewDocumentCommand{\gpunkte}{o m}{#2}
\newcommand{\gzuletzt}[1]{zuletzt #1}
% v0.7 (12.10.): Schülerfrage unter dem Stufenkopf; Ankreuzoption als Bild (Buchstabe und Kästchen darunter)
\newcommand{\gfrage}[1]{\par\nobreak\noindent\hspace*{\mbpfnr}{\small\itshape #1}\par\nobreak\addvspace{2pt}}
\newlength{\gbildbreite}\setlength{\gbildbreite}{0.32\linewidth}% Setzer setzt sie je Reihe (vier Bilder: 0,24)
\newcommand{\gbildoption}[3]{\begin{minipage}[t][#3][b]{\gbildbreite}\centering #2\par\smallskip\mbox{#1\,$\square$}\end{minipage}}
% Lineare (12.10.): Bild mit Buchstabe ohne Kästchen (Diagramm zuordnen); Aussage mit „wahr / falsch“ zum Ankreuzen
\newcommand{\gbildname}[3]{\begin{minipage}[t][#3][b]{\gbildbreite}\centering #2\par\smallskip\mbox{\textbf{#1}}\end{minipage}}
\newcommand{\gwahrfalsch}[2]{\par\noindent\hangindent2.2em\hangafter1\makebox[2.2em][l]{(#1)}#2\hfill\mbox{wahr $\square$\quad falsch $\square$}\par}
\makeatletter\setlength{\mbpfbe}{0mm}\makeatother
\begin{document}

\pfheftstil
\pflaufkopf[P10 \pkt{} mehr]{Lineare Funktionen}{PA10-9}\pfseitevon
\pfheftkopf{Lineare Funktionen}{mehr}
\pfstufenzeile{Punktprobe \pkt{} Gleichung aus Text \pkt{} Graph zeichnen \pkt{} Nullstelle gesucht \pkt{} Gleichung aus zwei Punkten}
% 2017-OS-K5b-m1 (Schwere I, rang 4)
\begin{pfaufg}{nach~P10\\’17~\pkt{}~5b}{1.}{}
Die Gerade $g$ hat die Gleichung $y = \tfrac{2}{3}x + 2$. Auf ihr liegt der Punkt $A(-12\,|\,y)$. Berechne seine y-Koordinate.
\karo{2}
\fusshilfe{\gfuss{1}{$-$6}}
\end{pfaufg}
% 2021-OS-K6a-m1 (Schwere I, rang 5)
\begin{pfaufg}{nach~P10\\’21~\pkt{}~6a}{2.}{}
Eine 36 cm hohe Kerze wird angezündet; nach 10 Minuten ist sie noch 33 cm hoch. Die Tabelle hält fest, wie hoch sie nach einigen Minuten noch ist. Trag die zwei fehlenden Werte in die Tabelle ein.
\par\smallskip\noindent\gpassend{{\small\renewcommand{\arraystretch}{1.2}\begin{tabular}{|l|c|c|c|c|c|c|c|}\hline Zeit in Minuten & 0 & 10 & 20 & 30 & 40 & \dots & 80\\\hline Höhe der Kerze in cm & \leerzelle & 33 & 30 & 27 & 24 & \dots & \leerzelle\\\hline\end{tabular}}}\par
\karo{3}
\fusshilfe{\gfuss{2}{0 min: 36 cm; 80 min: 12 cm}}
\end{pfaufg}
% 2021-OS-K6d-m1 (Schwere I, rang 6)
\begin{pfaufg}{nach~P10\\’21~\pkt{}~6d}{3.}{}
Eine 36 cm hohe Kerze wird angezündet. Die Tabelle hält fest, wie hoch sie nach einigen Minuten noch ist. Die Gleichung $y = -0{,}3x + 36$ beschreibt das Abbrennen der Kerze ($x$: Zeit in Minuten, $y$: Höhe in cm). Bestimme, nach wie vielen Minuten die Kerze ganz heruntergebrannt ist.
\par\smallskip\noindent\gpassend{{\small\renewcommand{\arraystretch}{1.2}\begin{tabular}{|l|c|c|c|c|c|c|c|}\hline Zeit in Minuten & 0 & 10 & 20 & 30 & 40 & \dots & 80\\\hline Höhe der Kerze in cm & \leerzelle & 33 & 30 & 27 & 24 & \dots & \leerzelle\\\hline\end{tabular}}}\par
\karo{2}
\fusshilfe{\gfuss{3}{120 min}}
\end{pfaufg}
% 2021-OS-K7a-m1 (Schwere I, rang 7)
\begin{pfaufg}{nach~P10\\’21~\pkt{}~7a}{4.}{}
Ordne jedem Sachverhalt die passende Gleichung zu. Sachverhalt 1: Eine Taxifahrt kostet 4 € Grundgebühr und 30 Cent je Kilometer. Sachverhalt 2: Ein Fitnessstudio verlangt im Monat 15 € Grundgebühr und 4 € für jeden gebuchten Kurs. Gleichungen: $y = 30x + 4$, $y = 4x + 15$, $y = 0{,}30x + 4$.
\par\smallskip\noindent\gpassend{{\small\renewcommand{\arraystretch}{1.3}\begin{tabular}{|l|c|c|c|}\hline & $y = 30x + 4$ & $y = 4x + 15$ & $y = 0{,}30x + 4$\\\hline Sachverhalt 1 & $\square$ & $\square$ & $\square$\\\hline Sachverhalt 2 & $\square$ & $\square$ & $\square$\\\hline\end{tabular}}}\par
\karo{3}
\fusshilfe{\gfuss{4}{4x + 15}}
\end{pfaufg}
% 2016-OS-K6a-m1 (Schwere I, rang 8)
\begin{pfaufg}{nach~P10\\’16~\pkt{}~6a}{5.}{}
Die 10b fährt mit dem Bus in eine Jugendherberge. Zwei Busfirmen machen ein Angebot: „Sonnenschein“ verlangt 150 € Grundpreis und 1,20 € je gefahrenem Kilometer, „Reiselust“ keinen Grundpreis und 1,80 € je Kilometer. Die Grafik zeigt für beide Angebote die Fahrtkosten in Abhängigkeit von der Strecke. Gib zu jedem der beiden Graphen die passende Busfirma an. Begründe für „Sonnenschein“, warum dein Graph passt.
\par\smallskip\noindent\gpassend{\begin{ksys}[xmin=0,xmax=350,ymin=0,ymax=700,xstep=25,ystep=50,karo=0.35,xlabel=Strecke in km,ylabel=Kosten in €]\gerade[300]{1.8}{0}{\mathrm{I}}\gerade[100]{1.2}{150}{\mathrm{II}}\end{ksys}}\par
\karo{3}
\fusshilfe{\gfuss{5}{I Reiselust, II Sonnenschein}}
\end{pfaufg}
% 2021-OS-K2a-m1 (Schwere I, rang 8)
\begin{pfaufg}{nach~P10\\’21~\pkt{}~2a}{6.}{}
Die lineare Funktion $f$ hat die Gleichung $y = 2x - 1$. Zeichne den Graphen von $f$ in das Koordinatensystem.
\par\smallskip\noindent\gpassend{\begin{ksys}[xmin=-4,xmax=4,ymin=-4,ymax=4,karo=0.5]\end{ksys}}\par
\karo{2}
\fusshilfe{\gfuss{6}{Gerade durch (0|$-$1) und (1|1)}}
\end{pfaufg}
% 2021-OS-K2b-m1 (Schwere I, rang 9)
\begin{pfaufg}{nach~P10\\’21~\pkt{}~2b}{7.}{}
Die lineare Funktion $f$ hat die Gleichung $y = 2x - 1$. Kreuze die beiden Aussagen an, die für $f$ stimmen:
\par\smallskip\noindent\pfkreuz{}{Der Punkt $Q(4\,|\,0)$ liegt auf dem Graphen.}\pfkreuz{}{Der Graph steigt monoton.}\pfkreuz{}{Der Graph trifft die y-Achse in $P(0\,|\,-1)$.}\pfkreuz{}{Der Graph geht durch den Ursprung.}
\fusshilfe{\gfuss{7}{Der Graph steigt monoton; Der Graph steigt monoton}}
\end{pfaufg}
% 2021-OS-K6c-m1 (Schwere II, rang 9)
\begin{pfaufg}{nach~P10\\’21~\pkt{}~6c}{8.}{}
Eine 36 cm hohe Kerze wird angezündet. Die Tabelle hält fest, wie hoch sie nach einigen Minuten noch ist. Das Abbrennen der Kerze beschreibt die Gleichung $y = -0{,}3x + 36$. Schreib zu $y$, $x$ und $36$, was sie hier bedeuten. Beispiel: $-0{,}3$ heißt, dass die Kerze in jeder Minute um 0,3 cm kürzer wird.
\par\smallskip\noindent\gpassend{{\small\renewcommand{\arraystretch}{1.4}\begin{tabular}{|c|l|}\hline Teil der Gleichung & Bedeutung\\\hline $y$ & \rule[-0.5ex]{45mm}{0.4pt}\\\hline $-0{,}3$ & In jeder Minute wird die Kerze 0,3 cm kürzer.\\\hline $x$ & \rule[-0.5ex]{45mm}{0.4pt}\\\hline $36$ & \rule[-0.5ex]{45mm}{0.4pt}\\\hline\end{tabular}}}\par
\karo{4}
\fusshilfe{\gfuss{8}{Höhe der Kerze in cm; Zeit in min ab Anzünden; Anfangshöhe in cm}}
\end{pfaufg}
% 2016-OS-K6b-m1 (Schwere II, rang 10)
\begin{pfaufg}{nach~P10\\’16~\pkt{}~6b}{9.}{}
Die 10b fährt mit dem Bus in eine Jugendherberge. Zwei Busfirmen machen ein Angebot: „Sonnenschein“ verlangt 150 € Grundpreis und 1,20 € je gefahrenem Kilometer, „Reiselust“ keinen Grundpreis und 1,80 € je Kilometer. Die Grafik zeigt für beide Angebote die Fahrtkosten in Abhängigkeit von der Strecke. Die Jugendherberge liegt 110 km entfernt; der Bus fährt hin und zurück. Gib an, welche Busfirma dafür billiger ist. Am Ort kommt noch ein Tagesausflug von 60 km hinzu. Untersuche, ob dieselbe Firma dann noch die billigere ist, und begründe.
\par\smallskip\noindent\gpassend{\begin{ksys}[xmin=0,xmax=350,ymin=0,ymax=700,xstep=25,ystep=50,karo=0.35,xlabel=Strecke in km,ylabel=Kosten in €]\gerade[300]{1.8}{0}{\mathrm{I}}\gerade[100]{1.2}{150}{\mathrm{II}}\end{ksys}}\par
\karo{3}
\fusshilfe{\gfuss{9}{Reiselust billiger: Sonnenschein 414 € gegen Reiselust 396; nein: bei 280 km Reiselust 504 €, Sonnenschein 486}}
\end{pfaufg}
% 2019-OS-K2a-m1 (Schwere I, rang 10)
\begin{pfaufg}{nach~P10\\’19~\pkt{}~2a}{10.}{}
In einem Koordinatensystem sind die Graphen $f$, $g$ und $h$ von drei linearen Funktionen gezeichnet. Gib die Nullstelle der Funktion $g$ an.
\par\smallskip\noindent\gpassend{\begin{ksys}[xmin=-5,xmax=5,ymin=-5,ymax=5,karo=0.5]\gerade[0.6]{-3}{3}{f}\gerade[1]{1}{3}{g}\gerade[3.5]{0}{-1}{h}\end{ksys}}\par
\karo{2}
\fusshilfe{\gfuss{10}{$-$3}}
\end{pfaufg}
% 2016-OS-K6c-m1 (Schwere II, rang 11)
\begin{pfaufg}{nach~P10\\’16~\pkt{}~6c}{11.}{}
Die 10b fährt mit dem Bus in eine Jugendherberge. Zwei Busfirmen machen ein Angebot: „Sonnenschein“ verlangt 150 € Grundpreis und 1,20 € je gefahrenem Kilometer, „Reiselust“ keinen Grundpreis und 1,80 € je Kilometer. Die Grafik zeigt für beide Angebote die Fahrtkosten in Abhängigkeit von der Strecke. Gib die Gleichung der Funktion an, deren Graph II ist (x: Strecke in km, y: Kosten in €).
\par\smallskip\noindent\gpassend{\begin{ksys}[xmin=0,xmax=350,ymin=0,ymax=700,xstep=25,ystep=50,karo=0.35,xlabel=Strecke in km,ylabel=Kosten in €]\gerade[300]{1.8}{0}{\mathrm{I}}\gerade[100]{1.2}{150}{\mathrm{II}}\end{ksys}}\par
\karo{2}
\fusshilfe{\gfuss{11}{1,2x + 150 (x in km, y in €)}}
\end{pfaufg}
% 2016-OS-B1c-m1 (Schwere I, rang 12)
\begin{pfaufg}{nach~P10\\’16~\pkt{}~1c}{12.}{}
Im Koordinatensystem sind zwei Geraden $f$ und $g$ gezeichnet. Kreuze an, welche der beiden fällt:
\par\smallskip\noindent\gpassend{\begin{ksys}[xmin=-2,xmax=4,ymin=-3,ymax=2,karo=0.5]\gerade[2.5]{-1}{1}{f}\gerade[-1]{0.75}{-1.5}{g}\end{ksys}}\par
\par\smallskip\noindent\pfkreuz{}{$f$}\pfkreuz{}{$g$.}
\fusshilfe{\gfuss{12}{f}}
\end{pfaufg}
% 2017-OS-K5a-m1 (Schwere II, rang 14)
\begin{pfaufg}{nach~P10\\’17~\pkt{}~5a}{13.}{}
Eine Gerade $g$ geht durch die Punkte $K(-6\,|\,-2)$ und $L(3\,|\,4)$. Zeichne ein Koordinatensystem und darin die Gerade $g$. Zeige, dass $y = \tfrac{2}{3}x + 2$ eine Gleichung von $g$ ist.
\par\smallskip\noindent\gpassend{\begin{tikzpicture}\draw[black!25,line width=0.3pt,step=0.5] (0,0) grid (12,7);\end{tikzpicture}}\par
\karo{4}
\fusshilfe{\gfuss{13}{Gerade durch K und L (y-Achsenabschnitt 2, Nullstelle $-$3); $\tfrac{2}{3}$x + 2}}
\end{pfaufg}
% 2015-OS-K4d-m1 (Schwere II, rang 15)
\begin{pfaufg}{nach~P10\\’15~\pkt{}~4d}{14.}{}
Lena springt mit dem Fallschirm aus einem Flugzeug. Ihre Höhe $h$ (in Metern) hängt von der Zeit $t$ (in Sekunden) ab. Nach 25 Sekunden freiem Fall öffnet sie in 1500 m Höhe den Schirm; 100 Sekunden später ist sie noch 900 m hoch. Während sie mit offenem Schirm gleitet, liegt der Graph auf einer Geraden. Stelle eine Gleichung dieser Geraden auf.
\par\smallskip\noindent\gpassend{\begin{ksys}[xmin=0,xmax=300,ymin=0,ymax=3000,xstep=25,ystep=250,karo=0.45,xlabel=$t$ in s,ylabel=$h$ in m]\funktionab{-2*\x^2+2750}{}{0}{25}\funktionab{-6*\x+1650}{}{25}{275}\end{ksys}}\par
\karo{2}
\fusshilfe{\gfuss{14}{$-$6t + 1650}}
\end{pfaufg}
\par\vfill\noindent{\scriptsize\color{mbgrau}Originale: P10 2017 \pkt{} 5b, P10 2021 \pkt{} 6a, P10 2021 \pkt{} 6d, P10 2021 \pkt{} 7a, P10 2016 \pkt{} 6a, P10 2021 \pkt{} 2a, P10 2021 \pkt{} 2b, P10 2021 \pkt{} 6c, P10 2016 \pkt{} 6b, P10 2019 \pkt{} 2a, P10 2016 \pkt{} 6c, P10 2016 \pkt{} 1c, P10 2017 \pkt{} 5a, P10 2015 \pkt{} 4d\par}
\end{document}